\chapter{De stelling van Pythagoras}\label{ch:g8:pythagoras}

De beroemdste stelling van de hele wiskunde legt een verband tussen de drie
zijden van een \omterm{def:g6:shapes:triangles}{rechthoekige driehoek}. Ermee worden lengten berekenbaar die geen
enkele meetlat rechtstreeks kan meten: diagonalen, afstanden, hoogten. Het is
ook onze eerste grote ontmoeting met het \omterm{ex:g1:subtraction:difference}{verschil} tussen een stelling en haar
omgekeerde.

\section{De stelling}

\begin{theorem}[Pythagoras]\label{thm:g8:pythagoras:direct}
Is een driehoek $ABC$ \omterm{def:g6:shapes:triangles}{rechthoekig} in $A$, dan geldt
\[
BC^2 = AB^2 + AC^2 :
\]
het kwadraat van de schuine zijde (de zijde tegenover de \omterm{def:g3:shapes:rightangle}{rechte hoek}) is de \omterm{def:g1:addition:def}{som}
van de kwadraten van de twee andere zijden.
\end{theorem}

\begin{proof}[Idee van het bewijs]
Vier exemplaren van de driehoek, samengelegd binnen een groot vierkant met zijde
$AB + AC$, laten een gekanteld vierkant op de schuine zijde onbedekt: het
vergelijken van de \omterm{def:g6:measure:area}{oppervlakten} geeft de gelijkheid. De uitgewerkte berekening is
\cref{exo:g8:pythagoras:11}; de stelling volgt ook uit het scalair \omterm{def:g2:mult:def}{product}, dat
in het bovenbouwvolume aan bod komt.
\end{proof}

\begin{omfigure}
\begin{tikzpicture}[scale=0.58]
  \coordinate (A) at (0,0);
  \coordinate (B) at (4,0);
  \coordinate (C) at (0,3);
  \draw[thick, fill=omDef!12] (A) -- (B) -- (C) -- cycle;
  \draw[thin] (0.35,0) -- (0.35,0.35) -- (0,0.35);
  \node[below left, font=\small] at (A) {$A$};
  \node[below right, font=\small] at (B) {$B$};
  \node[above left, font=\small] at (C) {$C$};
  % vierkanten op de zijden
  \draw[omDef, thick, fill=omDef!8] (0,0) rectangle (4,-4);
  \node[omDef, font=\small] at (2,-2) {$AB^2 = 16$};
  \draw[omMeth, thick, fill=omMeth!8] (0,0) rectangle (-3,3);
  \node[omMeth, font=\small] at (-1.5,1.5) {$AC^2 = 9$};
  \draw[omThm, thick, fill=omThm!8] (4,0) -- (7,4) -- (3,7) -- (0,3)
    -- cycle;
  \node[omThm, font=\small] at (3.5,3.5) {$BC^2 = 25$};
\end{tikzpicture}

{\small De stelling als uitspraak over \omterm{def:g6:measure:area}{oppervlakte}: het vierkant op de schuine
zijde ($25$) heeft precies dezelfde \omterm{def:g6:measure:area}{oppervlakte} als de twee andere vierkanten
samen ($16 + 9$). Hier zijn de zijden $3$, $4$, $5$.}
\end{omfigure}

\begin{method}[Een lengte berekenen]\label{met:g8:pythagoras:compute}
In een \omterm{def:g6:shapes:triangles}{rechthoekige driehoek} waarvan twee zijden bekend zijn:
\begin{enumerate}
  \item benoem de schuine zijde (tegenover de \omterm{def:g3:shapes:rightangle}{rechte hoek} --- altijd de langste
        zijde);
  \item schrijf de gelijkheid van \cref{thm:g8:pythagoras:direct} op met de
        bekende waarden;
  \item los op naar het ontbrekende kwadraat: \emph{tel} de kwadraten op als de
        schuine zijde onbekend is, \emph{trek} af als een rechthoekszijde
        onbekend is;
  \item neem de vierkantswortel, exact of met de rekenmachine, en ga na dat de
        schuine zijde er als langste zijde uit komt.
\end{enumerate}
\end{method}

\begin{example}[De schuine zijde vinden]\label{ex:g8:pythagoras:hypotenuse}
Een \omterm{def:g6:shapes:triangles}{rechthoekige driehoek} heeft rechthoekszijden $6$ en $8$. Voor de schuine
zijde $c$ geldt:
\[
c^2 = 6^2 + 8^2 = 36 + 64 = 100,
\qquad
c = \sqrt{100} = 10 .
\]
(Het getal $\sqrt{100}$, de \emph{vierkantswortel} van $100$, is het positieve
getal met kwadraat $100$; vierkantswortels krijgen een volledig hoofdstuk in
\cref{ch:g9:sqrt}.)
\end{example}

\begin{example}[Een rechthoekszijde vinden]\label{ex:g8:pythagoras:leg}
Een ladder van $5$ m staat tegen een muur, met zijn voet $1.4$ m van de muur.
Voor de bereikte hoogte $h$ geldt, met de ladder als schuine zijde,
\[
h^2 = 5^2 - 1.4^2 = 25 - 1.96 = 23.04,
\qquad
h = \sqrt{23.04} = 4.8 \text{ m}.
\]
Aftrekken, niet optellen: de onbekende is hier een \emph{rechthoekszijde}.
\end{example}

\section{De omgekeerde stelling}

\begin{theorem}[Omgekeerde stelling van Pythagoras]\label{thm:g8:pythagoras:converse}
Gelden in een driehoek $ABC$ voor de zijden $BC^2 = AB^2 + AC^2$, dan is de
driehoek \omterm{def:g6:shapes:triangles}{rechthoekig} in $A$.
\end{theorem}
\admitted

\begin{method}[Een rechte hoek toetsen]\label{met:g8:pythagoras:test}
Gegeven de drie zijden van een driehoek:
\begin{enumerate}
  \item bereken afzonderlijk het kwadraat van de langste zijde, en de \omterm{def:g1:addition:def}{som} van de
        kwadraten van de twee andere;
  \item zijn de twee uitkomsten gelijk, dan is de driehoek \omterm{def:g6:shapes:triangles}{rechthoekig}
        (omgekeerde stelling), en wel in het \omterm{def:g5:solids:def}{hoekpunt} tegenover de langste zijde;
  \item verschillen ze, dan is de driehoek niet \omterm{def:g6:shapes:triangles}{rechthoekig} --- want de stelling
        van Pythagoras zou dan gelijkheid opleggen.
\end{enumerate}
\end{method}

\begin{example}\label{ex:g8:pythagoras:test}
Zijden $5$, $12$, $13$: het kwadraat van de langste zijde is $13^2 = 169$; de \omterm{def:g1:addition:def}{som}
van de andere kwadraten is $5^2 + 12^2 = 25 + 144 = 169$. Gelijk: \omterm{def:g6:shapes:triangles}{rechthoekig}
(tegenover de zijde $13$).

Zijden $6$, $7$, $9$: $9^2 = 81$, maar $6^2 + 7^2 = 85 \neq 81$: geen
\omterm{def:g6:shapes:triangles}{rechthoekige driehoek} (hij is er ``bijna'' een --- de getallen beslissen, niet
het oog).
\end{example}

\begin{remark}[Stelling tegenover omgekeerde stelling]
De stelling en haar omgekeerde zeggen verschillende dingen: de stelling
\emph{gebruikt} een \omterm{def:g3:shapes:rightangle}{rechte hoek} om lengten te berekenen; de omgekeerde gebruikt
lengten om een \omterm{def:g3:shapes:rightangle}{rechte hoek} te \emph{bewijzen}. Bouwers gebruiken de omgekeerde al
duizenden jaren: een koord met knopen op $3$, $4$ en $5$ eenheden, strak
getrokken, geeft een volmaakt \omterm{def:g3:shapes:rightangle}{rechte hoek}.
\end{remark}

\begin{example}[Diagonaal van een vierkant]\label{ex:g8:pythagoras:diagonal}
De diagonaal van een vierkant met zijde $1$ is de schuine zijde van een
\omterm{def:g6:shapes:triangles}{rechthoekige driehoek} met rechthoekszijden $1$ en $1$:
\[
d^2 = 1^2 + 1^2 = 2 ,
\]
dus is $d = \sqrt2 \approx 1.414$ --- een getal dat geen \omterm{def:g6:fractions:def}{breuk} is, zoals het
bovenbouwvolume zal bewijzen. Voor een vierkant met zijde $c$ is de diagonaal
$c\sqrt2$.
\end{example}

\section{Oefeningen}

\begin{exercise}[$\star$]\label{exo:g8:pythagoras:1}
Bereken in elke \omterm{def:g6:shapes:triangles}{rechthoekige driehoek} de schuine zijde:
rechthoekszijden $9$ en $12$; \quad rechthoekszijden $5$ en $12$; \quad
rechthoekszijden $8$ en $15$.
\end{exercise}

\begin{exercise}[$\star$]\label{exo:g8:pythagoras:2}
Een \omterm{def:g6:shapes:triangles}{rechthoekige driehoek} heeft schuine zijde $25$ en één rechthoekszijde $7$.
Bereken de andere rechthoekszijde.
\end{exercise}

\begin{exercise}[$\star$]\label{exo:g8:pythagoras:3}
Bereken de diagonaal van een rechthoek met zijden $12$ cm en $9$ cm; en de
diagonaal van een vierkant met zijde $5$ cm (exacte waarde met een
vierkantswortel, daarna afgerond op de mm).
\end{exercise}

\begin{exercise}[$\star$]\label{exo:g8:pythagoras:4}
Welke driehoeken zijn \omterm{def:g6:shapes:triangles}{rechthoekig}? Verantwoord met
\cref{met:g8:pythagoras:test}:
\[
(20,\ 21,\ 29); \qquad
(7,\ 8,\ 11); \qquad
(1.5,\ 2,\ 2.5).
\]
\end{exercise}

\begin{exercise}[$\star$]\label{exo:g8:pythagoras:5}
Een poort wordt verstevigd met een schuine plank. De poort is $3$ m breed en
$1.6$ m hoog: hoe lang is de plank?
\end{exercise}

\begin{exercise}[$\star$]\label{exo:g8:pythagoras:6}
Een \omterm{def:g6:shapes:triangles}{gelijkbenige driehoek} heeft twee zijden van $10$ cm en een basis van $12$ cm.
Zijn hoogte verdeelt hem in twee \omterm{def:g6:shapes:triangles}{rechthoekige driehoeken}: bereken die hoogte, en
daarna de \omterm{def:g6:measure:area}{oppervlakte} van de driehoek.
\end{exercise}

\begin{exercise}[$\star\star$]\label{exo:g8:pythagoras:7}
Een televisiescherm is een $16{:}9$-rechthoek van $88.5$ cm bij $49.8$ cm.
Schermen worden verkocht op hun diagonaal, in inch ($1$ inch $= 2.54$ cm).
Bereken de diagonaal in cm, en daarna in inch (afgerond op de dichtstbijzijnde
inch).
\end{exercise}

\begin{exercise}[$\star\star$]\label{exo:g8:pythagoras:8}
Een boot vaart $24$ km naar het oosten en daarna $10$ km naar het noorden. Hoe
ver ligt hij dan in vogelvlucht van zijn vertrekpunt? Maak eerst een tekening.
\end{exercise}

\begin{exercise}[$\star\star$]\label{exo:g8:pythagoras:9}
Een ladder van $2.5$ m moet een vensterbank op $2.4$ m boven de grond bereiken.
Op welke afstand van de muur moet zijn voet staan? Past het antwoord bij het
veiligheidsadvies (voet op ongeveer een \omterm{def:g3:division:half}{kwart} van de ladderlengte van de muur)?
\end{exercise}

\begin{exercise}[$\star\star$]\label{exo:g8:pythagoras:10}
Zet in een assenstelsel de punten $A(1, 2)$ en $B(5, 5)$, en teken de
\omterm{def:g6:shapes:triangles}{rechthoekige driehoek} met rechthoekszijden \omterm{def:g6:lines:perp}{parallel} met de assen waarvan $[AB]$
de schuine zijde is. Bereken $AB$. (Deze constructie, voor twee willekeurige
punten, wordt de afstandsformule van de analytische meetkunde, in het
bovenbouwvolume.)
\end{exercise}

\begin{exercise}[$\star\star\star$]\label{exo:g8:pythagoras:11}
Vier exemplaren van een \omterm{def:g6:shapes:triangles}{rechthoekige driehoek} met rechthoekszijden $a$ en $b$ en
schuine zijde $c$ worden in de hoeken van een vierkant met zijde $a + b$ gelegd;
in het midden blijft een gekanteld vierkant met zijde $c$ onbedekt.
\begin{enumerate}
  \item Druk de \omterm{def:g6:measure:area}{oppervlakte} van het grote vierkant op twee manieren uit:
        rechtstreeks, en als (vier driehoeken) $+$ (gekanteld vierkant).
  \item Werk $(a + b)^2$ uit (zie \cref{thm:g8:equations:expand}) en leid daaruit
        $c^2 = a^2 + b^2$ af: je hebt de stelling van Pythagoras bewezen.
\end{enumerate}
\end{exercise}

\section{Opgave: pythagorese drietallen}

\begin{problem}[{Weekendopgave --- rechthoekige driehoeken met hele getallen,
en het merkwaardige product $(m^2 - n^2)^2 + (2mn)^2 = (m^2 + n^2)^2$ dat ze
allemaal voortbrengt}]
\label{pb:g8:pythagoras:1}
Een \emph{pythagorees drietal}\index{pythagorees drietal} is een drietal hele
getallen $(a, b, c)$ met
\[
a^2 + b^2 = c^2 ,
\]
zoals de $(3, 4, 5)$ van de bouwers. Een Babylonisch kleitablet bracht er
enorme in kaart --- $(4601,\ 4800,\ 6649)$ is er één van --- meer dan duizend
jaar vóór Pythagoras. Deze opgave jaagt eerst op drietallen en bouwt daarna de
machine die ze voortbrengt: een algebraïsche gelijkheid, rechtstreeks uit
\cref{ch:g8:equations}, in dienst van de meetkunde.

\medskip
\noindent\textbf{Deel I --- Op jacht naar drietallen.}
\begin{enumerate}
  \item Ga na dat $(5, 12, 13)$, $(8, 15, 17)$ en $(20, 21, 29)$ pythagorese
        drietallen zijn. Wat zegt \cref{thm:g8:pythagoras:converse} over
        driehoeken met deze zijden?
  \item Leg uit waarom een strak koord met twaalf gelijke stukken, uitgelegd als
        een driehoek met zijden van $3$, $4$ en $5$ stukken, bouwers een
        volmaakt \omterm{def:g3:shapes:rightangle}{rechte hoek} geeft.
  \item Toon aan: is $(a, b, c)$ een \omterm{pb:g8:pythagoras:1}{pythagorees drietal}, dan is
        $(ka, kb, kc)$ dat ook, voor elk heel getal $k \geq 1$. Leid uit
        $(3, 4, 5)$ drie nieuwe drietallen af.
  \item De \omterm{def:g5:division:multiple}{veelvouden} van $(3, 4, 5)$ zijn de drietallen $(3k, 4k, 5k)$. Toon
        aan dat $(5, 12, 13)$ daar \emph{niet} bij hoort.
  \item Leg uit waarom vragen 3 en 4 samen bewijzen dat het vermenigvuldigen van
        $(3, 4, 5)$ nooit elk \omterm{pb:g8:pythagoras:1}{pythagorees drietal} kan opleveren: er is een andere
        machine nodig.
\end{enumerate}

\medskip
\noindent\textbf{Deel II --- De machine.}
Neem twee hele getallen $m > n \geq 1$ en vorm de drie getallen
\[
a = m^2 - n^2,
\qquad
b = 2mn,
\qquad
c = m^2 + n^2 .
\]
\begin{enumerate}[resume]
  \item Werk met de merkwaardige \omterm{def:g2:mult:def}{producten} van \cref{pb:g8:equations:1} (schrijf
        $A = m^2$ en $B = n^2$) de kwadraten $a^2$ en $b^2$ uit, en bewijs dat
        \[
        \left(m^2 - n^2\right)^2 + (2mn)^2
        = \left(m^2 + n^2\right)^2 :
        \]
        de machine levert altijd een \omterm{pb:g8:pythagoras:1}{pythagorees drietal}.
  \item Laat de machine lopen op $(m, n) = (2, 1)$, $(3, 1)$, $(3, 2)$,
        $(4, 1)$ en $(4, 3)$, en zet de vijf drietallen in een tabel.
  \item Eén drietal in je tabel is een vermenigvuldigd exemplaar van een kleiner
        drietal. Welk, en van welk?
  \item Zoek de $(m, n)$ waarvoor de machine het drietal $(20, 21, 29)$ van
        vraag 1 oplevert (rechthoekszijden in willekeurige volgorde).
  \item Zoek $m$ en $n$ met $m^2 + n^2 = 100$, en leid daaruit een \omterm{pb:g8:pythagoras:1}{pythagorees
        drietal} af met schuine zijde $100$.
\end{enumerate}

\medskip
\noindent\textbf{Deel III --- Families, en een slot over pariteit.}
\begin{enumerate}[resume]
  \item Laat de machine lopen met $m = n + 1$. Toon aan dat de \omterm{def:g3:numbers:evenodd}{oneven}
        rechthoekszijde $2n + 1$ is, de \omterm{def:g3:numbers:evenodd}{even} rechthoekszijde $2n^2 + 2n$, en de
        schuine zijde $2n^2 + 2n + 1$ --- zodat de schuine zijde en de \omterm{def:g3:numbers:evenodd}{even}
        rechthoekszijde precies $1$ verschillen. Schrijf de drietallen uit voor
        $n = 1, 2, 3, 4$.
  \item Leid af dat \emph{elk} \omterm{def:g3:numbers:evenodd}{oneven getal} $3, 5, 7, 9, \dots$ een zijde is van
        een \omterm{def:g6:shapes:triangles}{rechthoekige driehoek} met hele getallen, en geef een drietal waarin
        $11$ voorkomt.
  \item Ook \omterm{def:g3:numbers:evenodd}{even} getallen lukken: laat de machine lopen met $n = 1$ en breng een
        drietal voort waarin $14$ voorkomt. Verkrijg daarna hetzelfde drietal een
        tweede keer, door een drietal uit je tabel van vraag 7 te
        vermenigvuldigen.
  \item Toon aan dat het kwadraat van een \omterm{def:g3:numbers:evenodd}{even getal} \omterm{def:g3:numbers:evenodd}{even} is en dat het kwadraat
        van een \omterm{def:g3:numbers:evenodd}{oneven getal} \omterm{def:g3:numbers:evenodd}{oneven} is. (Schrijf het getal als $2k$ of $2k + 1$
        en gebruik een merkwaardig \omterm{def:g2:mult:def}{product}.)
  \item Besluit het slot over pariteit: er is geen \omterm{pb:g8:pythagoras:1}{pythagorees drietal} waarvan
        alle drie de getallen \omterm{def:g3:numbers:evenodd}{oneven} zijn --- in elke \omterm{def:g6:shapes:triangles}{rechthoekige driehoek} met
        hele getallen is minstens één zijde \omterm{def:g3:numbers:evenodd}{even}.
\end{enumerate}
\end{problem}
